\section{AI Framework and Methodology}
\label{sec:methodologies}

\subsection{Data sources}

\red{The AI-enabled decision-support framework combines four hourly
information streams for Dozza.} Pedestrian data come from camera-level access
counts.  Mobile-network data\footnote{Data was bought from the Italian provider
TIM and are not publicly available} are provided as 15-minute aggregate
presence indicators and are summed over spatial tiles before being summarized
hourly.
\red{Weather data was obtained from the ERG5 gridded meteorological dataset
produced by ARPAE Emilia-Romagna \cite{ARPAE2023ERG5}. ERG5 provides hourly
variables, including temperature, precipitation, humidity, wind, solar
irradiance, leaf wetness, and reference evapotranspiration, derived through
spatial interpolation of observations from the regional meteorological station
network; this use is consistent with the high-resolution gridded climate
methodology documented by Antolini et al.~\cite{Antolini2016ERG5Climate}. The
local extract contains the two ERG5 grid cells used for the Dozza area, and the
modeling matrix retains both cell-level variables and cross-cell summaries such
as mean, minimum, maximum, and standard deviation.}

\red{Event features are derived from records downloaded locally by the project
scraper and then frozen before the cluster runs\footnote{Repository URL:
\url{https://github.com/FraResca/RetryDozza}.}. The
records cover Dozza and nearby relevant cities. They are transformed into
distance-band features, meaning active-event counts within fixed radii from
Dozza such as 10, 30, and 80 km; city and category indicators, meaning separate
active-event and intensity variables for places such as Dozza, Imola, Bologna,
province-capital cities, and event classes such as cultural, food-wine,
sport/motor, fair/congress, and music/show events; event-scale variables, where
lower values indicate small local events and higher values indicate larger
city-level festivals, fairs, concerts, or motor-sport events; event intensity,
defined as a distance-weighted and scale-weighted measure of active or imminent
event activity; and lagged or rolling event-activity features.}

The full source-level counts cover raw pedestrian rows, hourly pedestrian
timestamps, TIM input records, weather rows, and event-feature columns. In the
paper body, Table~\ref{tab:dataset-profile} focuses on the aligned modeling
matrix after the inner join. The joined period is the portion where pedestrian
targets, TIM aggregates, weather, and event features are simultaneously
available. Its profile shows a sparse, right-tailed pedestrian-count problem
paired with smoother TIM aggregates and a mostly spring-summer weather window.

\begin{table}
\caption{Descriptive profile of the aligned modeling dataset after joining pedestrian targets, TIM aggregates, weather variables, and event features.}
\label{tab:dataset-profile}
\centering
\scriptsize
\setlength{\tabcolsep}{3pt}
\begin{tabularx}{\textwidth}{lXX}
\toprule
Aspect & Value & Diagnostic \\
\midrule
Aligned temporal window & 2025-03-04 00:00 to 2025-09-30 23:00 & 2,735 observed hours over 5,064 calendar hours (54.0\% coverage) \\
Weekend representation & 30.4\% of aligned rows & Weekend hours are present but weekdays dominate the overlapping TIM period \\
Entrance-count scale & median 21, p90 92, max 587 & Hourly pedestrian counts are right-tailed, with occasional large peaks \\
Exit-count scale & median 27, p90 91, max 683 & Exit counts have a similar skewed scale to entrance counts \\
TIM total-presence scale & mean 5,462.3, p90 6,721.1 & TIM presences are smoother and larger than camera counts \\
TIM nationality mix & foreign share 2.1\% on average & Foreign presences are a small-scale target relative to Italian presences \\
TIM age composition & 51--60 share 22.2\%, 61+ share 39.6\% & Older age bands form the largest part of the Tc demographic aggregate \\
Weather conditions & mean temperature 18.5 C; rainy hours 9.5\% & The overlap period is mostly spring-summer and contains relatively few rainy hours \\
Event-signal density & 75.5\% of rows have at least one event feature active & Event variables are frequent but vary by distance, city, category, and intensity \\
Core-field missingness & 0.0\% average missingness & The inner join produces complete core target, TIM, weather, and event fields \\
\bottomrule
\end{tabularx}
\end{table}

Figure~\ref{fig:flow-target-distribution} shows the distributional shift between
the training and final test portions for the two pedestrian-flow targets. The
test split contains a higher upper range for exits and a slightly higher central
level for both flow targets. This shift explains why the evaluation uses a
chronological split and rolling validation rather than random cross-validation.

\begin{figure}
\centering
\includegraphics[width=0.86\textwidth]{flow_target_distribution_paper.pdf}
\caption{Pedestrian-flow target distributions in the chronological train and test splits. Outliers are hidden in the boxplot to keep the central range readable.}
\label{fig:flow-target-distribution}
\end{figure}

The pedestrian-flow targets are defined from the camera streams. The entrance
target ingressi\_borgo is the sum of entrances recorded by the Arcoribellino and
Piazzarocca Ingresso cameras. The exit target uscite\_borgo is the sum of exits
recorded by the Arcoribellino and Piazzarocca Uscita cameras. A timestamp is
marked complete only when the cameras needed for the target are present in the
same hour.

The TIM target families are derived from hourly mean-over-15-minute statistics.
The nationality track models Italian and foreign presences, while the age track
models six age-band presence indicators. The operational interpretation of the
TIM variables is shown in Table~\ref{tab:tim-dictionary}. The meanings are based
on the project data constraints and on public TIM City Forecast material
describing aggregated and anonymized mobile-network analytics for presence and
mobility monitoring \cite{TIMCityForecast,TIMAssisi}.

\begin{table}
\caption{Operational interpretation of the TIM fields used to derive presence, nationality, age, and mobility indicators.}
\label{tab:tim-dictionary}
\centering
\scriptsize
\begin{tabularx}{\textwidth}{lX}
\toprule
Fields & Meaning in this study \\
\midrule
event\_time; spatial fields & 15-minute timestamp and spatial identifiers for region, province, municipality, ACE area, and grid tile. \\
P, Ni, Ns & Total presences, Italian presences, and foreign presences; Ni + Ns = P in the raw data. \\
Tc, Tb & Italian consumer users with demographic detail and Italian business/non-consumer users; their sum approximates Ni. \\
Gm, Gf & Male and female aggregates within Tc; small differences may reflect rounding or unknown categories. \\
F1--F6 & Age bands within Tc: $<18$, 18--30, 31--40, 41--50, 51--60, and 61+. \\
Vp, Vr, Vi, Ve & Italian mobility categories: commuters, residents, intraregional visitors, and extraregional visitors. \\
mean15, sum15, max15 & Hourly mean, sum, and maximum over the 15-minute observations; sum15 is a person-quarter-hour aggregate. \\
\bottomrule
\end{tabularx}
\end{table}

\subsection{Preprocessing}

The preprocessing pipeline operates at hourly resolution. Pedestrian rows are
timestamped, grouped by hour, and aggregated into camera-level counts and
targets for the borough. TIM rows are converted from UTC to the Europe/Rome
analysis timezone, aggregated across tiles at 15-minute resolution, and then
summarized by hour using mean, sum, and maximum statistics. Weather measurements
from the local ERG5 grid are joined by hourly timestamp. Event records
are normalized into hourly active-count, intensity, scale, distance-band,
city-level, category-level, lagged, and rolling features. The final modeling
dataset is an inner join between complete pedestrian targets and TIM hourly
features, with weather and event covariates aligned on the same hourly
timestamp.

The inner join intentionally sacrifices the full pedestrian time span in order to
retain only hours where the main sources overlap. The pedestrian archive spans
from March 2022 to September 2025, while the joined modeling window is
restricted by TIM availability in 2025. The resulting testable dataset contains
2,735 hourly rows and 175 columns. This design is appropriate for cross-source
forecasting, but conclusions are limited to the observed TIM-overlap period.

\subsection{Forecasting setup and feature selection}

Experiments are run in forecast mode at 1, 3, 6, 12, and 24 hour horizons. The
main detailed tables report the one-hour horizon, while the multi-horizon
summaries and figures are used to assess robustness. The final evaluation uses a
chronological split with 2,164 training rows and 541 test rows. The train period
runs from 2025-03-04 18:00:00 to 2025-09-05 07:00:00; the test period runs from
2025-09-05 08:00:00 to 2025-09-30 22:00:00. Rolling validation summaries over
three folds are computed for every track and horizon.

\begin{table}
\caption{Experimental configuration used for the multi-horizon forecasting runs.}
\label{tab:setup}
\centering
\begin{tabularx}{\textwidth}{lX}
\toprule
Item & Value \\
\midrule
Forecast mode & forecast \\
Prediction horizons & 1, 3, 6, 12, and 24 hours \\
Primary reported horizon & 1 hour, with multi-horizon robustness tables \\
Train split & 2,164 rows, 2025-03-04 18:00:00 to 2025-09-05 07:00:00 \\
Test split & 541 rows, 2025-09-05 08:00:00 to 2025-09-30 22:00:00 \\
Lags & 1, 2, 24, and 168 hours \\
Rolling windows & 3, 6, and 24 hours \\
Execution nodes & CPU jobs ran on cnode04.hpc.prv and cnode06.hpc.prv, both in the cnode0[3-6] class: 2 x AMD EPYC 9454 48-Core CPUs and 512 GB DDR5 RAM per node \\
Slurm resources & Model jobs requested 4 CPU threads; 16 GB RAM for flow/nationality jobs and 24 GB RAM for age jobs \\
Top-k feature selection & Rolling validation over $k \in \{15,25,30,40,60\}$ for every track and horizon \\
Models & Persistence baselines, Ridge/log-Ridge, Poisson, Tweedie, Random Forest, Extra Trees, HGB, XGBoost, LightGBM, and two-stage Ridge \\
Interpretability & Permutation importance, SHAP, leave-one-group-out ablation, leave-one-feature-out ablation, and event-feature ablation summaries \\
\bottomrule
\end{tabularx}
\end{table}

Feature selection is leakage-aware. In forecast mode, same-hour variables from
the target family are excluded when they would reveal the future target, while
lagged and rolling historical variables are allowed. At the one-hour horizon,
the flow and nationality tasks start from 1,237 candidate features after leakage
exclusion and retain 265 and 264 features, respectively, after missingness and
collinearity filtering. The age task starts from 1,225 candidates and retains
264. The final top-k value is selected by rolling validation for every track and
horizon: at one hour, the selected values are 25 for flow and 15 for both
nationality and age.

The leakage rule is stricter than direct target removal. In forecast mode, all
unlagged TIM and pedestrian variables are excluded from feature models, because
they correspond to same-hour information at the forecast target time. Only
lagged and rolling versions of TIM and pedestrian variables are eligible. This
also prevents same-hour algebraic relationships among TIM fields, such as total
presence, nationality fields, consumer/business fields, and age-band fields,
from revealing the target. Calendar, weather, and event variables at the
forecast origin are retained, and quality or completeness flags are always
removed.

\begin{table}
\caption{Rolling-validation top-k selection by horizon and target family.}
\label{tab:top-k}
\centering
\scriptsize
\begin{tabular}{rlrrrr}
\toprule
Horizon & Track & Selected $k$ & Targets & Mean validation MAE & Mean validation $R^2$ \\
\midrule
1 h & flow & 25 & 2 & 16.65 & 0.672 \\
1 h & nationality & 15 & 2 & 56.36 & 0.876 \\
1 h & age & 15 & 6 & 20.72 & 0.802 \\
3 h & flow & 15 & 2 & 17.55 & 0.657 \\
3 h & nationality & 15 & 2 & 56.85 & 0.872 \\
3 h & age & 15 & 6 & 20.89 & 0.798 \\
6 h & flow & 15 & 2 & 17.64 & 0.645 \\
6 h & nationality & 15 & 2 & 56.42 & 0.874 \\
6 h & age & 15 & 6 & 20.77 & 0.801 \\
12 h & flow & 15 & 2 & 17.54 & 0.650 \\
12 h & nationality & 15 & 2 & 57.07 & 0.876 \\
12 h & age & 15 & 6 & 20.92 & 0.801 \\
24 h & flow & 15 & 2 & 17.44 & 0.656 \\
24 h & nationality & 15 & 2 & 56.42 & 0.882 \\
24 h & age & 15 & 6 & 20.76 & 0.807 \\
\bottomrule
\end{tabular}
\end{table}

\subsection{Models, metrics, and interpretation}

The model roster contains simple baselines and feature-based learners. Baselines
include mean and median dummy predictors, last-hour persistence, same-hour
previous-day/week predictors, and 24/168-hour rolling-mean predictors.
Feature-based models include ridge, log-ridge, Poisson, Tweedie, Random Forest,
Extra Trees, histogram gradient boosting, XGBoost, LightGBM, and a two-stage
ridge model. The tree-ensemble methods follow standard tabular-learning practice
\cite{Geurts2006ExtraTrees,Chen2016XGBoost,Ke2017LightGBM}.

Models are evaluated using MAE, RMSE, $R^2$, SMAPE, WAPE, MASE, peak-hour
classification metrics, training time, and inference time per row. WAPE is
emphasized when comparing targets with different scales, while MASE provides a
scale-normalized comparison against naive dynamics. Interpretation uses
selected-feature group counts, leave-one-group-out ablation,
leave-one-feature-out ablation, permutation importance, SHAP summaries, and
specific summaries for event-feature selection and event-feature ablation.
These artifacts are used as diagnostics of predictive dependence, not as causal
estimates. Leave-one-group-out ablation re-trains the selected feature model
after removing an entire feature family. Leave-one-feature-out ablation is more
local: for each target it re-trains the lowest-MAE feature-based model after
removing one selected feature at a time, without re-running hyperparameter
search.
Timing values are wall-clock
measurements collected with the same Python timing calls inside the CPU Slurm
jobs described in Table~\ref{tab:setup}; they support internal model comparison
within this pipeline rather than hardware-independent benchmarking.
